\subsection{Background}

Skin cancer is among the most common cancers worldwide, and the prognosis for a
patient depends heavily on how early the lesion is found. Melanoma in
particular is highly treatable when caught early but becomes far more dangerous
once it spreads. Dermoscopy a non-invasive imaging technique that removes
surface glare and reveals the sub-surface structure of a lesion has become a
standard tool for this kind of early assessment. Reading dermoscopic images
well, however, takes years of training, and even experienced dermatologists do
not always agree with one another. This gap between the volume of lesions that
need screening and the number of specialists available to read them is what
motivates automated, image-based diagnostic support.

Deep learning has reshaped this area over the past decade. Convolutional neural
networks (CNNs) trained on large dermoscopic collections reached
dermatologist-level performance on melanoma recognition~\cite{esteva2017dermatologist},
and the annual ISIC challenges have since turned skin lesion classification into
a well-defined benchmark task. The ISIC-2019 dataset, which we use throughout
this work, contains 25{,}331 labelled dermoscopic images spanning eight
diagnostic categories: melanoma (MEL), melanocytic nevus (NV), basal cell
carcinoma (BCC), actinic keratosis (AK), benign keratosis (BKL),
dermatofibroma (DF), vascular lesion (VASC), and squamous cell carcinoma (SCC).

\subsection{Rationale of the Study or Motivation}

Two practical problems keep most published skin-lesion models from being used in
an everyday clinic. The first is cost. The strongest results on ISIC-2019 come
from large ensembles of heavy networks~\cite{isic2019leaderboard,gessert2020ensembles};
these are accurate but impractical to run on the modest hardware found in a
typical dermatology practice or on a portable device next to the patient. The
second is that image-only models ignore information a dermatologist always has
in front of them. The age of the patient, their sex, and the part of the body
where the lesion sits all shift the prior probability of each diagnosis a
scaly patch on the face of an elderly patient is read very differently from the
same patch on a young person's back. This clinical context is recorded in the
ISIC-2019 metadata, yet it is routinely discarded by classifiers that look only
at pixels.

Our study is driven by the belief that these two problems should be solved
together rather than separately. We argue first that a carefully designed
lightweight hybrid can match or beat much larger image-only baselines on
balanced accuracy while staying small enough to deploy, and second that the
clinical metadata a dermatologist already has on hand can be folded in at
negligible cost as additional diagnostic context, handled so that missing
fields are never imputed. We
also recognise that not every deployment has the same constraints, so we deliver
two models from one design philosophy: a genuinely compact model for
constrained settings and a higher-capacity model for situations where balanced
accuracy matters more than footprint.

\subsection{Problem Statement}

ISIC-2019 is severely imbalanced. The most common class, nevus, has roughly
12{,}875 images, while dermatofibroma and vascular lesions have only about 239
and 253 respectively a ratio of more than fifty to one. A model trained
naively on this data learns to predict the common classes and quietly ignores
the rare ones, which is the opposite of what matters clinically, since several
of the rare classes are malignant. Reporting overall accuracy hides this failure
entirely. The problem we address is therefore to build a model that performs
well \emph{across all eight classes at once} measured by balanced
multi-class accuracy (BMA), the mean of the per-class recalls while
respecting a tight parameter and compute budget and while making use of the
clinical metadata that accompanies each image, including the large fraction of
records where some of that metadata is missing.

\subsection{Objective}

The objectives of this work are:

\begin{enumerate}
  \item To design a lightweight hybrid CNN Transformer classifier for
        ISIC-2019 that stays under 6M parameters and 1 GMAC and that fuses
        clinical metadata through a cross-attention head capable of handling
        missing fields without imputation.
  \item To design a second, higher-capacity model in the same family a
        dual-backbone network with learned attention fusion and a
        Kolmogorov Arnold Network (KAN) classifier that trades footprint
        for a higher balanced accuracy.
  \item To train both models, together with five established lightweight
        baselines, under one identical loss, augmentation, and evaluation
        pipeline so that the comparison is fair.
  \item To quantify the contribution of each design choice (the CNN stem, the
        transformer trunk, the metadata head, and the dual-backbone fusion)
        through controlled ablations, and to report efficiency alongside
        accuracy.
\end{enumerate}

